% ===== CV TEMPLATE: Hardware Systems Engineering =====
% Target roles: embedded systems engineer, hardware development engineer,
%   electronics engineer, automation/control systems engineer, R&D hardware engineer
%   at large industrial companies (ABB, Atlas Copco, Ericsson, Axis, etc.)
% Emphasis: breadth of hardware engineering disciplines, system integration
%   ownership, embedded/control experience, proven ability to take complex
%   electromechanical systems from design to delivery.
% Claude Code instruction: copy to cv/cv.tex. Adjust the Objective to mirror
%   the specific role title and emphasis. For embedded/firmware-heavy roles,
%   expand the embedded/control bullet; for broader hardware roles keep as-is.
%   Do not alter facts or dates.
\documentclass[a4paper,11pt]{article}
\usepackage[utf8]{inputenc}
\usepackage[T1]{fontenc}
\usepackage[margin=1in]{geometry}
\usepackage{setspace}
\usepackage{hyperref}
\usepackage{enumitem}
\setstretch{1.0}
\pagestyle{empty}

\begin{document}

%====================
% Header
%====================
\begin{center}
    {\LARGE \textbf{Bo Yuan}} \\
    \vspace{0.15cm}
    +46 769\,612\,787 \enspace $\cdot$ \enspace
    \href{mailto:yuanbo330lab@gmail.com}{yuanbo330lab@gmail.com} \enspace $\cdot$ \enspace
    \href{mailto:boyua@kth.se}{boyua@kth.se} \\
    Stockholm / Uppsala, Sweden
\end{center}

\vspace{0.3cm}

%====================
% Objective
%====================
\section*{Objective}

Hardware systems engineer with a strong background in electronics design,
embedded and PLC-based control systems, electromechanical integration, and
full-lifecycle hardware development. Five years as CTO owning the design,
integration, and productisation of a complex high-power electromechanical system
from first principles to small-series delivery. Currently completing an MSc in
Electromagnetics and Plasma Engineering at KTH. Seeking hardware development or
systems engineering roles where broad technical ownership and cross-disciplinary
integration are valued.

%====================
% Education
%====================
\section*{Education}

\noindent\textbf{KTH Royal Institute of Technology}, Stockholm, Sweden
\hfill 09/2024 -- 07/2026 (expected) \\
\textit{M.Sc.\ in Electromagnetics, Fusion and Space Engineering --- Plasma Track} \\
Relevant coursework: Electromagnetic Waves, Plasma Physics, Microwave Modelling,
Vacuum Technology, Plasma--Surface Interaction, Fusion Systems

\medskip
\noindent\textbf{Xi'an University of Technology}, China \hfill 09/2013 -- 07/2018 \\
\textit{B.Eng.\ in Electrical Engineering and Automation}

%====================
% Technical Skills
%====================
\section*{Technical Skills}

\begin{itemize}[noitemsep, topsep=2pt, leftmargin=*]
    \item \textbf{Electronics and PCB Design:}
          Altium Designer ($\sim$10\,yr); multi-layer mixed-signal PCB design,
          schematic capture, layout, DFM review; coordination with EMS/SMT
          manufacturers; prototype hand-assembly, inspection, and bring-up

    \item \textbf{Embedded and Control Systems:}
          Siemens S7-1200 PLC (process control logic, I/O handling, signal
          conditioning); custom PCB-based control architecture replacing
          traditional control cabinets; basic C++ and Python for system
          tooling and data acquisition; MCU-based systems (prior experience)

    \item \textbf{Mechanical and Electromechanical Integration:}
          SolidWorks ($\sim$10\,yr); system-level mechanical design,
          electromechanical assembly, tolerance management, CNC fixture design,
          high-vacuum component design and sealing

    \item \textbf{RF and High-Power Systems:}
          High-power microwave systems (2.45~GHz, 10~kW); waveguide structures,
          impedance matching, vacuum-compatible RF components; high-voltage
          power supply integration

    \item \textbf{Simulation and Modelling:}
          COMSOL Multiphysics ($\sim$3\,yr), ANSYS HFSS ($\sim$2\,yr),
          Keysight ADS ($\sim$1\,yr); Python for data analysis and engineering scripts

    \item \textbf{System Integration and Verification:}
          Subsystem interface definition, module-level unit testing,
          vacuum leak testing, system-level acceptance testing,
          failure analysis and troubleshooting across hardware disciplines
\end{itemize}

%====================
% Professional Experience
%====================
\section*{Professional Experience}

\noindent\textbf{Sichuan 330 Semiconductor Ltd.}, China
\hfill 07/2019 -- 06/2024 \\
\textit{Chief Technology Officer / Co-Founder}

\noindent Owned the full development lifecycle of a high-power MPCVD
(Microwave Plasma Chemical Vapor Deposition) system --- a tightly coupled
electromechanical platform integrating high-power RF, high-vacuum,
thermal management, embedded control, and precision mechanics ---
from first-principles design through prototyping and small-series production,
delivering $\sim$20 complete systems. Company subsequently acquired.

\medskip
\noindent\textit{Electronics and Control:}
\begin{itemize}[noitemsep, topsep=2pt, leftmargin=*]
    \item Defined PCB-based electrical integration strategy to replace
          traditional control cabinets; reduced wiring complexity and
          improved electrical reliability across all builds
    \item Designed custom PCBs in Altium Designer for power distribution,
          signal conditioning, and sensor interfacing; coordinated
          external SMT manufacturing and managed prototype bring-up
    \item Implemented Siemens S7-1200 PLC for process control: interlock
          logic, gas flow, power ramp sequencing, and alarm handling;
          defined control architecture and I/O allocation
    \item Integrated high-voltage power supplies, magnetron microwave source,
          cooling loops, gas handling, and vacuum instrumentation into a
          unified system with defined electrical and mechanical interfaces
\end{itemize}

\medskip
\noindent\textit{System Architecture and Hardware Development:}
\begin{itemize}[noitemsep, topsep=2pt, leftmargin=*]
    \item Defined modular system architecture decoupling RF, vacuum, thermal,
          and control subsystems; enabled parallel development and independent
          module-level testing before system integration
    \item Drove DFM decisions throughout: consolidated multi-part assemblies
          into CNC-machined monolithic components, reducing tolerance
          stack-up, assembly steps, and field failure risk
    \item Designed complete high-vacuum subsystem: chamber geometry, sealing
          strategy, leak detection procedures, and dedicated vacuum leak test
          fixtures for production quality control
    \item Led the R\&D team (7 engineers) and directed development across
          mechanical, electrical, RF, and software disciplines simultaneously
\end{itemize}

\medskip
\noindent\textit{Selected Project Contributions (MIST Satellite --- KTH, 2025):}
\begin{itemize}[noitemsep, topsep=2pt, leftmargin=*]
    \item Served as mechanical group coordinator in KTH's MIST student
          satellite project; led debugging and troubleshooting across
          mechanical subsystems and contributed to functional subsystem tests
\end{itemize}

\medskip
\noindent\textbf{Chengdu Kechuang Institute of Science and Culture}, China
\hfill 03/2019 -- 06/2019 \\
\textit{Research Engineer}

\noindent Assessed technical feasibility of MPCVD-based diamond synthesis;
successfully raised 2\,M CNY in seed funding for Sichuan 330 Semiconductor Ltd.

%====================
% Patents
%====================
\section*{Selected Patents}

\begin{itemize}[noitemsep, topsep=2pt, leftmargin=*]
    \item \textbf{High-power TE--TEM microwave mode converter}
          (CN212162036U, 2020)
    \item \textbf{Waveguide sealing window} (CN212392359U, 2021)
    \item \textbf{Microwave sliding short with choking structure}
          (CN212162044U, 2020)
    \item \textbf{Large-area compact MPCVD device} (CN212560433U, 2021)
\end{itemize}

%====================
% Languages
%====================
\section*{Languages}

Chinese (native) \enspace $\cdot$ \enspace English (fluent, professional working level)

\end{document}
